% titel: Blatt 0, A3: Summennotation

% Teil I: Werkzeuge, kein Lösungsweg.

% Grundidee: Das Summenzeichen ist eine Kompression von ``ausschreiben``.
% Der Index ist ein Laufplatzhalter; das Umschreiben verlangt, dass erster
% Term und letzter Term stimmen.

Das Summenzeichen $\sum$ ist keine neue Rechenoperation, sondern
\emph{Kurzschrift} für ``addiere viele gleich gebaute Terme``. Wer die
Notation versteht, versteht auch, warum sie so praktisch ist: man gibt
nur an, wie der allgemeine Summand aussieht und \emph{wo} er beginnt und
endet.

\begin{werkzeugkasten}[Summenzeichen: Aufbau]
Eine Summe über die Laufvariable $i$ von $i=m$ bis $i=n$ schreibt man
\begin{equation}
  \sum_{i=m}^{n} a_i = a_m + a_{m+1} + \dots + a_n .
\end{equation}
\textbf{Vier Bausteine:} das Zeichen, der Laufindex $i$ (die ``Frage
welches Glied``), der Startwert $m$ und der Endwert $n$. Die Summanden
$i=m,m+1,\dots,n$ werden eingefüllt und aufsummiert.

Der Laufindex $i$ ist \emph{stumm}: Er ist nur ein Platzhalter innerhalb
der Summe und darf umbenannt werden, ohne den Wert zu ändern,
\begin{equation}
  \sum_{i=m}^{n} a_i = \sum_{k=m}^{n} a_k .
\end{equation}
Anschaulich: ``das Glied Nummer $i$`` und ``das Glied Nummer $k$``
beschreiben dieselbe Liste, nur mit anderem Namensschild.
\end{werkzeugkasten}

\begin{werkzeugkasten}[Warum man beim Umbenennen den \emph{Startwert} anpasst]
Will man einen Summanden umschreiben (z.\,B. $2k+1$ statt $k$ als
Laufvariable), kommt es nur auf zwei Terme an:
\begin{itemize}
  \item \textbf{Startwert:} Der erste Summand der neuen Schreibweise muss
        mit dem ersten der alten übereinstimmen. Beispiel: Ungerade Zahlen
        $1,3,5,\dots$ starten mit $1$; mit der Form $2k+1$ ist das
        $k=0$.
  \item \textbf{Endwert:} Der letzte Term gibt den passenden Endwert vor.
\end{itemize}
Das Umschreiben ist also eine ``Passprobe an den Rändern``: Stimmen
Anfang und Ende, stimmt die ganze Summe.
\end{werkzeugkasten}

\begin{achtung}[Typische Fehler]
\begin{itemize}
  \item Start- und Endwert des Index falsch ablesen. ``$1+2+3+\dots+10$``
        ist $\sum_{i=1}^{10}i$, nicht bis $9$ oder $11$.
  \item Die Summanden nicht ausrechnen --- es ist nur die \emph{andere
        Schreibweise} verlangt.
  \item Exponenten richtig interpretieren (z.\,B. $2^2,3^2,\dots,k^2$,
        nicht die Ziffern ``zusammensetzen``).
  \item Beim Ausformulieren von $\sum_{i=a}^{b}f(i)$ alle Terme für
        $i=a,a+1,\dots,b$ konkret hinschreiben --- besonders den ersten
        und den letzten, weil man sie für die Ränder-Probe braucht.
\end{itemize}
\end{achtung}
